{
    \setlength{\parindent}{0em}
    \par {\zihao{4}\bfseries 一、题目：\uline{\hfill \Title \hfill}}
    \par {\zihao{4}\bfseries 二、指导教师对开题报告的具体要求：}
    
}

本科毕业设计是对学生综合能力与学术素养的全面检验，
要求学生围绕毕业设计主题，针对核心技术广泛查阅近三年国内外重要期刊或会议论文，
其中英文文献不少于20篇、中文文献不少于2篇，总参考文献不少于30篇，
并翻译一篇英文论文，完成系统性文献综述，以准确把握研究现状与发展趋势。
随后，学生需撰写开题报告（参考学校格式），阐明研究背景与意义、国内外研究现状及现有方案比较，
明确研究内容与预期目标，提出可行的研究方案与技术路线，制定详细进度安排，经指导教师审阅后参加开题答辩。
在设计过程中，学生应按计划推进研究，撰写毕业设计论文，确保论文结构完整、字数达标、格式规范，
并通过查重检测。答辩时，学生以PPT形式清晰陈述研究成果，回答答辩委员会提问，
答辩委员会结合论文质量与现场表现综合评定成绩。整个过程强调学术诚信与AI规范使用，严禁抄袭，
所有材料均需指导教师签字确认，以确保毕业设计高质量完成。

\mbox{} \vfill

\signature{指导教师（签名）}
